
\begin{proposition}[Unicyclic quasi-bundle supports with all torsors retained]\label{cdp:st:C.2.UN.62}
Let $S$ be a specified compact connected SNC surface with smooth
embedded components. Each component is either a nonlinear primary
or secondary Hopf surface of algebraic dimension zero, or is
equipped with an actual map
\[
                         q_v:T_v\longrightarrow\PP^1
\tag{C.QT.1}\label{cdp:eq:C.QT.1}
\]
with connected fibres, every fibre a multiple of a smooth elliptic
curve, fixed elliptic Jacobian, and identity monodromy on the
integral first cohomology of the regular fibre. Every boundary
curve on $T_v$ is a fibre support of this specified map. The
elliptic gluing maps, the surfaces and their complete analytic
torsors are part of the data.

The following geometric test is necessary and sufficient for
$S$ to occur as a full reduced fibre of a smooth-total
degeneration of connected complex two-tori.
\begin{enumerate}
\item The actual boundary normal lines are dual, and the incidence
graph has exactly one independent cycle. Nonlinear Hopf vertices
are leaves.
\item Every multiple fibre of every map \eqref{cdp:eq:C.QT.1} is marked by
a boundary curve. Let $m_e$ be the exact order of its normal
line, also the normal order at the other endpoint. At a
nonlinear endpoint this is its resonance.
\item Set $h_v=0$ on the cycle and set $h_v$ equal to the sum
of the $m_e$ along its path from the cycle. There is an integer
$R>0$ for which every nonlinear leaf has height $R$ and every
elliptic quasi-bundle vertex has height strictly smaller than
$R$. With no nonlinear leaves one may choose any
$R>\max_v h_v$.
\item Apply the exact inverse-logarithmic and product-cap
pruning maps of \ref{cdp:st:C.2.UN.60}, retaining their full integer
framings and all torsor identifications. The resulting
marked diagonal Hopf or degree-zero ruled core cycle has
integral complement-lattice product $M=I$.
\end{enumerate}
All pruning operations in this test exist under its preceding
requirements. Their marked result is independent of the order
of disjoint branch operations. Every support passing the test
has a realization over $\PP^1$ with smooth compact total space,
retaining every original component, analytic torsor and
specified elliptic gluing map. The only other singular support
is the complex conjugate of its pruned core. Every reduced
fibre support has Euler number zero.

In particular the statement includes elliptic quasi-bundles
of Kodaira dimension one and coherent Euler characteristic
zero, together with the genus-zero quasi-bundle components
of Kodaira dimension zero and the Hopf and ruled components
in the same explicitly specified class.
\end{proposition}

\begin{proof}
\par\medskip\noindent\textbf{The actual root chart and its cohomology maps.}
Let $mJ$ be a fibre of \eqref{cdp:eq:C.QT.1}. The normalized local root
$t=s^m$ is finite unramified over the surface tube, by
the exact chart $s=wx$, $w^m=u$ for $t=ux^m$.
Its generic fibre is connected. Its finite Stein factor
over the root disc has degree one, so the central
elliptic cover is connected as well. The fixed Jacobian
and identity linear monodromy make the root family a
translation torsor under the fixed elliptic curve $F$.
The sequence $0\to\Lambda_F\to\OO\to F\to0$ and
$H^1(\Delta,\OO)=H^2(\Delta,\Lambda_F)=0$ trivialize
this actual torsor. The finite translation gauge in
(C.HT.4) therefore gives the surface tube exactly as
\[
 (F\times\Delta_s)/\mu_m,
 \qquad(z,s)\longmapsto(z+\tau,\zeta_ms),
 \qquad\operatorname{ord}(\tau)=m.
\tag{C.QT.2}\label{cdp:eq:C.QT.2}
\]
The last order is exact because the central action is
free. Its associated normal character is primitive.
Thus $N_{J/T_v}$ has exact order $m$, and its connected
normal-root cover is precisely $F\to J$.

The fixed cohomology frame $H^1(F,\OO_F)$ on this
root chart is unchanged by the finite translations.
Taking invariant functions gives the ordinary
coordinate $s^m$. The regular frames glue by the
specified identity monodromy and by translations.
Consequently the actual direct images and their
cohomology are
\[
 q_{v*}\OO_{T_v}=\OO_{\PP^1},\qquad
 R^1q_{v*}\OO_{T_v}=\OO_{\PP^1},\qquad
 (h^0,h^1,h^2)(T_v,\OO)=(1,1,0).
\tag{C.QT.3}\label{cdp:eq:C.QT.3}
\]
At a marked fibre $J$, restriction to $J$ followed
by the finite pullback $F\to J$ is evaluation on this
same cohomology frame. The resulting map
$H^1(T_v,\OO)\to H^1(F,\OO_F)$ is an isomorphism.
In particular the original restriction to $J$ is
nonzero. The commuting square is the actual root
chart \eqref{cdp:eq:C.QT.2}; no identification of different
cohomology groups by dimensions alone is used.

A nonlinear Hopf leaf has the same values $(1,1,0)$
and the nonzero restriction map proved in \ref{cdp:st:C.2.UN.54}.
Duality of boundary normal lines identifies their
exact orders and their connected normal-root covers.
It thus identifies the endpoint Jacobians by actual
finite-cover maps, including the secondary character
of a nonlinear leaf from \ref{cdp:st:C.2.UN.50}.

\par\medskip\noindent\textbf{Graph, residues and the canonical minimum.}
Distinct fibre supports on each $T_v$ are disjoint,
and a nonlinear Hopf has only its one elliptic curve.
Hence the specified SNC surface has no triple points.
For an actual degeneration its full fibre is reduced,
so the meridian quotient of \ref{cdp:st:C.2.UN.44} makes $\pi_1(S)$
abelian. The actual map onto the incidence graph
then makes that graph's free fundamental group
abelian. It has at most one independent cycle.

The normalization sequence is exactly
\[
 0\longrightarrow\OO_S\longrightarrow
   \bigoplus_v\OO_{T_v}\longrightarrow
   \bigoplus_e\OO_{J_e}\longrightarrow0,
\tag{C.QT.4}\label{cdp:eq:C.QT.4}
\]
with nonlinear Hopf terms included in the vertex
sum. The induced $H^1$ matrix is a weighted
incidence matrix whose every endpoint coefficient
is nonzero, by \eqref{cdp:eq:C.QT.3}. A tree gives a surjective
matrix by leaf elimination, and therefore
$H^2(S,\OO_S)=0$. Flatness over the smooth curve
and upper semicontinuity from a torus instead
give $h^2(S,\OO_S)\ge1$. Thus the graph is
unicyclic. Its spanning-tree minor has rank
$\#V-1$; the remaining row must vanish on that
minor's kernel. This is exactly product-one
holonomy for the actual endpoint restriction
maps, and
\[
                    (h^0,h^1,h^2)(S,\OO_S)=(1,2,1).
\tag{C.QT.5}\label{cdp:eq:C.QT.5}
\]

Serre duality identifies the transpose incidence
map with the component residue obstruction maps.
Its kernel has zero residues on every attached
tree, by leaf elimination, and nonzero residues
on every cycle edge. Since every vertex has
$H^0(K_{T_v})=0$, a component section with all
residues zero is zero. The unique global section
of $\omega_S$ is consequently nonzero exactly
on the cycle components. Divide the actual
canonical evaluation by its minimum full-fibre
power. Its nonzero reduction is this section.
Its minimum vertical order $k_0$ is therefore
attained exactly on the cycle, as in the explicit
evaluation comparison of \ref{cdp:st:C.2.UN.54}.

\par\medskip\noindent\textbf{The exact differential and height calculation.}
On a vertex \eqref{cdp:eq:C.QT.1}, invariant fibre differentials
give actual meromorphic two-forms
$q_v^*\theta\wedge dz$. Translation charts
change $dz$ only by a multiple of the base
differential, which disappears in this wedge.
The forms therefore glue. Every meromorphic
two-form with vertical divisor has this form:
its ratio to one reference form is constant
on the regular compact fibres and descends;
\eqref{cdp:eq:C.QT.2} makes its descended meromorphic
function invariant under $s\mapsto\zeta_ms$,
so it is meromorphic in $s^m$ across a multiple
fibre. The local order is exactly
\[
 \operatorname{ord}_{J}(q_v^*\theta\wedge dz)
             =m\operatorname{ord}_{q_v(J)}\theta+m-1.
\tag{C.QT.6}\label{cdp:eq:C.QT.6}
\]

Restriction of the canonical evaluation divided
by $t^{k_v}$ gives the actual meromorphic
component form with divisor
\[
           \sum_{e=(v,w)}(k_w-k_v-1)J_e.
\tag{C.QT.7}\label{cdp:eq:C.QT.7}
\]
An unmarked multiple fibre would have order zero
in \eqref{cdp:eq:C.QT.7}, and \eqref{cdp:eq:C.QT.6} would require
$m\operatorname{ord}\theta+m-1=0$, hence $m=1$.
At a marked fibre these same equations give
\[
 d_{v,e}:=\frac{k_w-k_v}{m_e}
       =\operatorname{ord}_{q_v(J_e)}\theta+1\in\Z,
 \qquad\sum_{e\ni v}d_{v,e}=\operatorname{val}(v)-2.
\tag{C.QT.8}\label{cdp:eq:C.QT.8}
\]
The last equality is the degree $-2$ of an
actual rational differential on $\PP^1$.
Opposite orientations negate $d$ because their
normal-root covers have the same exact order.

At a terminal elliptic vertex \eqref{cdp:eq:C.QT.8} gives
height difference $m_e$ from its parent. A
nonlinear terminal vertex gives precisely the
same difference by its invariant axis form
$dz\wedge dw/w^{m_e+1}$ and the exact secondary
normal comparison in \ref{cdp:st:C.2.UN.50}. Induct inward along
each tree. If the $c$ child terms are one,
\eqref{cdp:eq:C.QT.8} gives $-d_{\rm parent}+c=c-1$,
so the parent term is also one. The cycle
terms are zero by the canonical minimum
already identified. Hence $k_v-k_0$ is exactly
the stated path sum $h_v$.

This proof also gives the actual base differentials:
on a cycle vertex it is $c\,dz/z$, with the two
cycle points $0,\infty$; on a tree vertex it is
$c\,dz$, with its parent point at infinity.
The other marked points have order zero.
No Kodaira-dimension assumption was used in
this calculation.

\par\medskip\noindent\textbf{The common integral period.}
Dual normal-root charts give on an edge
$xy=t$, $q_v=x^{m_e}$ and $q_w=y^{m_e}$,
with the same connected elliptic cover and
the actual normal-product trivialization.
The level $x^{m_e}=c$ corresponds to
$y^{m_e}=t^{m_e}/c$. Varying its nonzero
normal level gives the oriented isotopy
between the two regular-fibre two-tori.
The product-preserving tubular comparison
transfers this equality to the actual
nearby torus, as in (C.HT.10).
At a nonlinear leaf the unimodular boundary
map of \ref{cdp:st:C.2.UN.50} identifies its real core with
this same integral class. All vertices
therefore give one class $\gamma$.

The finite-root restriction square in
\eqref{cdp:eq:C.QT.3} also shows that the coherent
cycle holonomy in \eqref{cdp:eq:C.QT.5} is the linear
action of the actual product of elliptic
root-cover isomorphisms. Its value one
forces that product to be a translation.
This retains its possible translation
and all Jacobian identifications.

The period of the evaluation divided by
$t^{k_0}$ is holomorphic. On a regular
holomorphic fibre curve of $T_v$ its
leading two-form restricts to zero, so
\[
                  P(t)=O(t^{h_v+1}).
\tag{C.QT.9}\label{cdp:eq:C.QT.9}
\]
At a nonlinear leaf of secondary order
$n$ and contraction parameter $\beta$,
the exact primitive core computation is
\[
 P(t)=t^{h_v}
       \left(\frac{\pm2\pi i c\beta^{-m_e}}n+O(t)\right),
                         \qquad c\ne0.
\tag{C.QT.10}\label{cdp:eq:C.QT.10}
\]
Thus every nonlinear leaf has the one
finite order $R$ of this actual common
period, and \eqref{cdp:eq:C.QT.9} makes every elliptic
vertex's height smaller. Without such a
leaf this condition places no further
restriction: an integer above the finite
set of heights exists. This proves all
the height requirements for actual fibres.

\par\medskip\noindent\textbf{Pruning with the full torsors.}
The pruning proof of \ref{cdp:st:C.2.UN.60} applies to the
maps \eqref{cdp:eq:C.QT.1} themselves. We specify the
comparisons that retain the present
broader surface class. Delete a terminal
nonlinear Hopf by the inverse chart
(C.PR.6), with its full primitive period
basis and integer $n$. After the children
of an elliptic vertex have been removed,
its actual modified map still has fixed
Jacobian and identity monodromy over
$\PP^1$; each inverse operation replaced
one multiple fibre by the displayed
ordinary product tube and changed no
other transition. Its only possible
multiple fibre is now its parent mark.

It is therefore the actual $F$-torsor
on $\PP^1[(\infty,m)]$ classified by
(C.PR.1)--(C.PR.5). The exact logarithm
sequence and $H^1(\OO)=H^2(\OO)=0$
identify the entire torsor with its
two root-line classes. A nonzero
vector gives the literal weighted
Hopf quotient, including $n$ and
the primitive fibre direction;
a zero vector forces $m=1$ and
the product ruled cap. This conclusion
concerns the actual modified surface,
even when its original unmodified
surface had Kodaira dimension one.
Remove that pair by its corresponding
map and retain the entire inverse
record.

At the end the two-mark root-stack
cover (C.PR.7) and its actual line
pushout identify each core vertex
with a marked diagonal Hopf surface
or the torsion ruled bundle (C.PR.8).
These are exact maps on the original
transition atlases. Disjoint branch
operations commute on those atlases,
so their marked result has the
independence stated in \ref{cdp:st:C.2.UN.60}.

For an actual degeneration, every
terminal complement and its inverse
product cap have the identical
boundary homomorphism (C.PR.9),
retaining both elliptic generators
and killing exactly the specified
primitive meridian. The van Kampen
pushouts are identical on the
unchanged exterior. Their iteration
identifies the core plumbing group
with the original $\Z^4$. The
injected-lattice semidirect product
for that core consequently has
$M=I$. This proves necessity of
the last test by actual group maps.

\par\medskip\noindent\textbf{Compact restoration of every prescribed component.}
Conversely suppose the test passes.
Use the compact marked core of \ref{cdp:st:C.2.UN.57}
with the power $V_{t^R}$ and the
unchanged logarithmic coefficient
of (C.PR.11). Its necks remain
$xy=t$ and its opposite filling
is the conjugate pruned core.
The compact elliptic subgroup of
its common open group is the
actual regular-fibre group, and
(C.PR.12) makes its transverse
period determinant a nonzero
constant times $t^R$.

Reverse the entire pruning record.
At a vertex of height $h$, the
local quotient (C.PR.13) has
displacement of order $R-h$.
The recorded primitive basis is
implemented by the displayed
unimodular linear period map
following (C.PR.13). At a Hopf
step, the weight-$(m,1)$ insertion
with $\OO(nE)$ has exactly the
free finite chart (C.PR.14).
The new child's coordinate
$U=Z/W^m$ gives the actual
displacement $\delta(t)/t^m$,
as in (C.PR.15). At a product
cap the ordinary blowup has
the exact displacement
$\delta(t)/t$, as in (C.PR.16).
The height test makes all these
maps holomorphic and gives
precisely the prescribed diagonal
or nonlinear central type.

All operations use the entire
recorded inverse transitions.
In particular an original
properly elliptic vertex is
restored by reversing the
successive changes of its
root lines and ordinary tubes;
its complete analytic torsor,
all multiple fibres and all
marked points reappear, rather
than only its degree or normal
characters. The strict-parent
map is exactly (C.PR.6) in the
opposite direction. The same
argument applies to every
other original quasi-bundle.
Fibre translations and the
recorded linear maps retain
each specified elliptic gluing.

The norm and affine-base bounds
(C.PR.17)--(C.PR.18) prove the
actual local properness and
the old-cluster comparison for
each replacement. Their constants
depend on the local displacement
order and contraction, not on the
Kodaira dimension of the surface
being restored. The finitely many
disjoint replacements therefore
give the same Hausdorff compact
gluing as \ref{cdp:st:C.2.UN.60}. The weighted
free coordinates have $t=xW$;
ordinary cap charts have $t=xy$.
Every final component consequently
has full multiplicity one.

Each component \eqref{cdp:eq:C.QT.1} has Euler
number zero by additivity over
its elliptic fibre supports;
each nonlinear Hopf also has
Euler number zero. Every double
curve is elliptic and there are
no triple points. Both special
supports thus have Euler zero,
as do the general tori. This
proves the complete criterion
and the compact assertion.

\par\medskip\noindent\textbf{An explicit properly elliptic component.}
Fix $F=\C^*/\beta^{\Z}$, $0<|\beta|<1$.
Start with the two-component
ruled core $F\times\PP^1$
glued along its two marked
sections by the identity.
Its lattice product is the
identity. On one component
choose four distinct ordinary
fibre points away from those
sections. Perform four recorded
weight-three, secondary-order-one
insertions, with the common
period order $R=3$.
They add four nonlinear primary
Hopf leaves $H_3(\beta)$.

The modified central component
has an explicit independent
description. Put
\[
 \mathcal B=\PP^1[(p_1,3),\ldots,(p_4,3)],\qquad
 L=\OO_{\mathcal B}(D_1+\cdots+D_4),
 \qquad 3D_i=p_i.
\tag{C.QT.11}\label{cdp:eq:C.QT.11}
\]
It is the smooth coarse surface
of $L^\times/\beta^{\Z}$.
Each inertia acts by a primitive
fibre translation, so its coarse
surface is smooth and its four
fibres are exactly triple smooth
elliptic curves. In a root chart
$p=s^3$ its line coordinate is
$y=sW$, where $W$ is the old
product coordinate. The inverse
chart (C.PR.6) with $m=3,n=1$
has precisely $W=y/s$.
Thus this given line pushout
is the component produced by
the four insertions, including
its full torsor.
Explicitly the positive divisor
line has frames $e_{\rm in}=s^{-1}e_{\rm out}$.
The equality of the actual vector
$W e_{\rm out}=y e_{\rm in}$
is precisely $y=sW$; its canonical
section has local coefficient $s$.
This fixes the sign of each $D_i$
in \eqref{cdp:eq:C.QT.11}.

For any map of the class
\eqref{cdp:eq:C.QT.1}, the actual invariant
fibre differential and the
root chart give the plurigenus
formula
\[
 h^0(T_v,\omega_{T_v}^{\otimes a})
 =\max\left\{0,\ 1-2a+
       \sum_i\left\lfloor
                    a(1-1/m_i)\right\rfloor\right\}.
\tag{C.QT.12}\label{cdp:eq:C.QT.12}
\]
Indeed a pluricanonical section
is a base meromorphic
$a$-differential times the
$a$th invariant fibre factor;
holomorphy at a root chart
permits precisely the displayed
pole order. Descent and its
inverse are those actual
invariant-coordinate maps.
For \eqref{cdp:eq:C.QT.11} this gives
\[
 h^0(T,\omega_T^{\otimes a})
    =\max\{0,1-2a+4\lfloor2a/3\rfloor\},
 \qquad \omega_T^{\otimes3}\simeq q^*\OO_{\PP^1}(2).
\tag{C.QT.13}\label{cdp:eq:C.QT.13}
\]
In particular $p_g(T)=0$,
$q(T)=1$, $\chi(\OO_T)=0$
and $h^0(T,\omega_T^{\otimes3b})=2b+1$.
Thus this is an actual properly
elliptic component. The complete
compact reduced fibre consists
of it, the unchanged ruled core
component, and the four Hopf
leaves, with every gluing given
by the preceding maps.
The root-line maps, the incidence
matrix, the root-coordinate pole
bounds and these exact plurigenera
are checked in
\path{code/D_quasibundle_trees.py}
and its saved output.
\end{proof}


